\chapter{Mathematical Preliminaries}\label{c1}
These notes follow the chapters of \cite{so1989} as a general route towards 
understanding methods of quantum chemistry. The book was published in 1989 and
a lot of resources easily available now were not within a student's reach four
decades ago. I hope to make plenty of digressions as I study these chapters.

\section{The variation method}
We want to find an approximate solution to the time-independent Schr\"{o}dinger
equation
\begin{equation}\label{c1e1}
\opr{H}\ket{\Phi} = \mathcal{E}\ket{\Phi}
\end{equation}
Suppose that $\opr{H}$ has a discrete spectrum of eigenvalues $\{\mathcal{E}_n\}$
with $n = 0, 1, \ldots$. Let $\{\ket{\Phi_n}\}$ be the corresponding eigenkets.
Since $\opr{H}$ is hermitian, the eigenvalues are real and the eigenkets can be
assumed to be orthonormal. Therefore, we also have,
\begin{equation}\label{c2e2}
\braket{\Phi_a|\opr{H}|\Phi_b} = E_a\delta_{ab}.
\end{equation}
We can also assume that $\{\ket{\Phi_a}\}$ form a complete set so that any other
ket $\ket{\tilde{\Phi}}$ obeying the same boundary conditions can be expressed
as 
\begin{equation}\label{c2e3}
\ket{\tilde{\Phi}} = \sum_a\ket{\Phi_a}\braket{\Phi_a|\tilde{\Phi}}.
\end{equation}

In almost all cases, except the simplest system, $\{\ket{\Phi_a}\}$ are unknown.
But their existence can be assumed following the mathematical properties of
$\opr{H}$. Therefore, \ref{c2e3} is true in principle. If $\ket{\tilde{\Phi}}$
is an approximate solution to \eqref{c1e1} then we will show that
\begin{equation}\label{c2e4}
\braket{\tilde{\Phi}|\opr{H}|\tilde{\Phi}} \ge \mathcal{E}_0.
\end{equation}
That is, the true ground state energy of the system is a lower bound to the
expectation value of $\opr{H}$ in the state $\ket{\tilde{\Phi}}$. Therefore, in
any method to find an approximate solution to \eqref{c1e1} we choose the ket
that minimises the expectation value of $\opr{H}$ in it. Typically, we consider
a family of \emph{trial kets} that depend on a set of parameters. We choose the
parameters that minimise the expectation value of the hamiltonian in the state.
Before illustrating this method for a few simple problems, we will prove the
inequality \eqref{c2e4}. Let
\begin{equation}\label{c2e5}
\ket{\tilde{\Phi}} = \sum_a c_a\ket{\Phi_a}
\end{equation}
so that,
\[
\opr{H}\ket{\tilde{\Phi}} = \sum_a c_aE_a\ket{\Phi_a}
\]
and
\begin{equation}\label{c2e6}
\braket{\tilde{\Phi}|\opr{H}|\tilde{\Phi}} = \sum_{a,b}c_b^\ast c_a\mathcal{E}_a
\braket{\Phi_b|\Phi_a} \sum_a |c_a|^\mathcal{E}_a.
\end{equation}
Since $\ket{\tilde{\Phi}}$ can be assumed to be normalised. 
\[
\sum_a |c_a^2| = 1.
\]
If the eigenstates are labelled such that $\mathcal{E}_0 \le \mathcal{E}_1 \le
\ldots$ then
\[
\braket{\tilde{\Phi}|\opr{H}|\tilde{\Phi}} \ge \mathcal{E}_0\sum_a|c_a|^2 =
\mathcal{E}_0.
\]

Examples:
\begin{enumerate}
\item This is exercise 1.18 in the book \cite{so1989}. We will first solve
the hamiltonian analytically. Let
\[
\left(-\frac{1}{2}\frac{d^2}{dx^2} - \delta(x)\right)\psi(x)=\mathcal{E}\psi(x).
\]
In the region $x \ne 0$ the solution is
\[
\psi(x)=\begin{dcases}Ae^{-\kappa x} & x > 0 \\Be^{\kappa x} & x < 0\end{dcases}
\]
where $A$ and $B$ are constants. Since $\psi$ must be continuous at $x = 0$,
we must have $A = B$. This gives us the energy eigenvalue
\[
\mathcal{E} = -\frac{\kappa^2}{2}.
\]
If we integrate the Schr\"{o}dinger equation over the region $[-a, a]$ then
we have
\[
-\frac{1}{2}(\psi^\op(a) - \psi^\op(-a)) - \psi(0) = \mathcal{E}
\int_{-a}^a\psi(x)dx.
\]
so that
\[
-\frac{1}{2}(-\kappa A e^{-\kappa a} - \kappa A e^{+\kappa a}) = 
\psi(0) + \mathcal{E}\int_{-a}^a\psi(x)dx.
\]
In the limit $a \rightarrow 0$, since $\psi$ is continuous, the integral 
vanishes and this equation gives
\[
\kappa = 1
\]
Thus the true energy eigenvalue is $-1/2$.

We will now use the variational method to get it. Let $\tilde{\Phi} = Ne^{-ax^2}$
be a trial function. The normalisation condition gives 
\[
N^2 = \sqrt{\frac{2a}{\pi}}.
\]
We next evaluate
\[
\opr{H}\tilde{\Phi} = aN(1 - 2ax^2)e^{-ax^2} - \delta(x)Ne^{-ax^2}.
\]
so that
\[
\braket{\opr{H}}(a) = \int_{-\infty}^\infty dx\left(aN^2(1 - 2ax^2)e^{-2ax^2} - 
\delta(x)N^2 e^{-2ax^2}\right) = \frac{a}{2} - \sqrt{\frac{2a}{\pi}}.
\]
This function has an extremum at $a_0 = 2/\pi$ so that
\[
\braket{\opr{H}}(a_0) = -\frac{1}{\pi} \approx -0.3183 > -0.5.
\]
We used equation \eqref{a1e8} to evaluate the integrals in this example.

\item This is exercise 1.19 in the book \cite{so1989}. The hamiltonian for the
hydrogen atom is
\[
\opr{H} = \left(-\frac{1}{2}\nabla^2 - \frac{1}{r}\right).
\]
Once again we try the gaussian trial function
\[
\ket{\tilde{\Phi}} = Ne^{-ar^2}.
\]
This is a three dimensional wave-function and therefore the normalisation
condition is
\[
\int_0^\infty\int_0^{2\pi}\int_0^\pi N^2 e^{-2ar^2}r^2\sin\theta drd\theta d\phi
= 4\pi N^2 \frac{1}{8a}\sqrt{\frac{\pi}{2a}} = 1
\]
so that
\[
N^2 = \sqrt{\frac{2a}{\pi}}
\]
Now,
\[
\opr{H}\ket{\tilde{\Phi}} = aN(3 - 2ar^2)e^{-ar^2} - N\frac{e^{-ar^2}}{r}
\]
so that
\[
\braket{\opr{H}}(a) = \frac{3a}{2} - 2\sqrt{\frac{2a}{\pi}}.
\]
Its minimum is as 
\[
a_0 = \frac{8}{9\pi}
\]
so that the minimum expectation value is
\[
\braket{\opr{H}}(a) = -\frac{4}{3\pi} \approx -0.4244,
\]
while the actual ground state energy is $-0.5$ hartree.


\end{enumerate}

\section{Standard integrals}
The most basic Gaussian integral is
\begin{equation}\label{a1e1}
\int_{-\infty}^\infty e^{-x^2}dx = \sqrt{\pi}.
\end{equation}
To evaluate it, start with 
\[
I^2 = \int_{-\infty}^\infty\int_{-\infty}^\infty e^{-x^2-y^2}dxdy
\]
Transforming to polar coordinates
\[
I^2 = \int_0^\infty\int_0^{2\pi} e^{-r^2}rdrd\theta = 
2\pi\int_0^\infty e^{-r^2}rdr = \pi
\]
from which equation \eqref{a1e1} follows immediately. From \eqref{a1e1} we 
have, using a quick transformation of variables
\begin{equation}\label{a1e2}
\int_{-\infty}^\infty e^{-ax^2}dx = \sqrt{\frac{\pi}{a}}.
\end{equation}
The integral
\[
\int_{-\infty}^\infty e^{-ax^2+bx}dx
\]
is evaluated by completing the square
\[
-ax^2+bx = -\left(ax^2 - 2\sqrt{a}x\frac{b}{2\sqrt{a}} + \frac{b^2}{4a}\right)
+ \frac{b^2}{4a} = -\left(ax - \frac{b}{2\sqrt{a}}\right)^2 + \frac{b^2}{4a}
\]
we have
\[
\int_{-\infty}^\infty e^{-ax^2+bx}dx = \exp\left(\frac{b^2}{4a}\right)
\int_{-\infty}^\infty\exp\left[-\left(ax-\frac{b}{2\sqrt{a}}\right)^2\right]dx
\]
so that
\begin{equation}\label{a1e3}
\int_{-\infty}^\infty e^{-ax^2+bx}dx = 
\exp\left(\frac{b^2}{4a}\right)\sqrt{\frac{\pi}{a}}.
\end{equation}
A trivial generalisation of this result is
\begin{equation}\label{a1e4}
\int_{-\infty}^\infty e^{-ax^2+bx+c}dx = \exp\left(\frac{b^2}{4a} + c\right)
\sqrt{\frac{\pi}{a}}.
\end{equation}
The integral \eqref{a1e1} is related to the gamma function. For, using the 
transformation $t = x^2$,
\[
\int_{-\infty}^\infty e^{-x^2}dx = 2\int_{0}^\infty e^{-x^2}dx =
\int_{0}^\infty e^{-t}t^{-1/2}dt = \Gamma\left(\frac{1}{2}\right)
\]
where
\begin{equation}\label{a1e5}
\Gamma(z) = \int_{0}^\infty t^{z-1}e^{-t}dt.
\end{equation}
Thus we have the frequently used result
\begin{equation}\label{a1e6}
\Gamma\left(\frac{1}{2}\right) = \sqrt{\pi}.
\end{equation}
Integrating \eqref{a1e5} by parts once we get
\begin{equation}\label{a1e7}
\Gamma(z) = (z - 1)\Gamma(z - 1).
\end{equation}

We now start with equation \eqref{a1e2} and differentiate it with respect to 
$a$ to get
\[
\int_{-\infty}^\infty -x^2e^{-ax^2}dx = -\frac{1}{2}\frac{\sqrt{\pi}}{a^{3/2}}
\]
or
\[
\int_{-\infty}^\infty x^2e^{-ax^2}dx = 
\frac{1}{2}\Gamma\left(\frac{1}{2}\right)\frac{1}{a^{3/2}} = 
\frac{\Gamma(3/2)}{a^{3/2}}.
\]
Integrating it once more with respect to $a$ you get
\[
\int_{-\infty}^\infty x^4e^{-ax^2}dx = \frac{\Gamma(5/2)}{a^{5/2}}.
\]
Integrating it $n$ times, we get the result
\begin{equation}\label{a1e8}
\int_{-\infty}^\infty x^{2n}e^{-ax^2}dx = 
\frac{\Gamma(\frac{2n+1}{2})}{a^{(2n+1)/2}}.
\end{equation}
and hence
\begin{equation}\label{a1e9}
\int_{0}^\infty x^{2n}e^{-ax^2}dx = 
\frac{1}{2}\frac{\Gamma(\frac{2n+1}{2})}{a^{(2n+1)/2}}.
\end{equation}
It is clear that
\begin{equation}\label{a1e10}
\int_{-\infty}^\infty x^{2n+1}e^{-ax^2}dx = 0.
\end{equation}
However, we can the integral with limits $0$ to $\infty$ will be non-zero. To 
get their value, we start with
\[
\int_0^\infty x e^{-ax^2}dx = \frac{1}{2a}
\]
and differentiate it $n$ times to get
\begin{equation}\label{a1e11}
\int_0^\infty x^{2n+1} e^{-ax^2}dx = \frac{n!}{2a^{n+1}}.
\end{equation}
By putting $2n+1=k$ in equation \eqref{a1e9} and $2n=k$ in equation 
\eqref{a1e11} we get the single formula
\begin{equation}\label{a1e12}
\int_0^\infty x^{k} e^{-ax^2}dx = 
\frac{\Gamma\left(\frac{k+1}{2}\right)}{2a^{(k+1)/2}}.
\end{equation}

\section{Constrained optimisation}
Consider the problem of finding the maximum of a function $f:\mathbb{R}^2
\rightarrow \mathbb{R}$ subject to the constraint $g(\vec{r}) = c$ where
$g:\mathbb{R}^3 \rightarrow \mathbb{R}$. If there was no constraint, we
would have equated $\grad f = 0$. However the constraint requires us to traverse
the surface $z = f(x, y)$ along the curve $g(x, y, z) = c$. In general, 
travelling along the curve $g(x, y, z) = c$ involves crossing the level curves
of $f$. However, at the extrema, the level curves of $f$ are tangential to $g$
or that the gradients of the two functions are parallel to each other. Therefore,
there exists a constant $\lambda$ such that $\grad f = \lambda\grad g$. 
Equivalently, the condition for constraind optimisation is
\begin{equation}\label{a1e1}
\grad (f - \lambda g) = 0.
\end{equation}
The constant $\lambda$ need not be calculated and is therefore called an
undetermined multiplier of the constraint.
